\section{Reproducibility Manifest}\label{app:manifest}
A complete study manifest should record: implementation hash and local diff;
Python/compiler/platform and native API/backend; full effective environment and
map; world seed policy and actual seeds; algorithm/network/observation/action
contracts per species; population/food/body/reproductive/mutation settings;
optimizer/update/trial budgets; champion/early-stop selection and source models;
recording cadence; independent replication count; primary endpoints; exclusions;
and analysis code/uncertainty estimator. Relative experiment paths are useful
locally but are not a substitute for an archived dataset identifier.

\subsection{Illustrative commands}
The following commands describe workflows, not the settings of an experiment
already reported in Supplementary Section~\ref{sec:observations}. Run them from the repository
root, in its isolated Python environment. Inspect the JSON for any explicit
initial-model configuration before calling a run fresh.
\begin{lstlisting}
# Fresh initialization under an incrementing, reproducible seed schedule:
.venv/bin/python -m ecosystem.cli environments/ppo.json \
  --runs 10 --seed 100 --brain prey=fresh --brain predator=fresh --progress

# Use archived compatible species/sex models in a new experiment:
.venv/bin/python -m ecosystem.cli environments/ppo.json \
  --brains-from results/PREVIOUS_EXPERIMENT \
  --brain prey=lifespan --brain predator=kills --runs 10 --progress

# Continue the exact saved experiment state:
.venv/bin/python -m ecosystem.cli \
  --resume results/EXPERIMENT/checkpoints/final.bin.xz

# Start the independent saved-run viewer:
./run.sh
\end{lstlisting}
With multiple runs, default chaining can use retained champions. To restart each
world from its original models, add \code{--no-chain-brains}. A completed
experiment or completed plateau stop remains completed on exact resume; use a
new warm-start experiment when the goal is additional independent episodes.
\code{--seed} selects the deterministic increment policy unless explicitly
changed by the available seed-policy controls; save the effective settings.

\section{Engineering Evidence Inventory}\label{app:evidence}
\begin{table}[ht]
\centering\small
\caption{Repository sources for the transcribed engineering evidence. Historical
report measurements should be rerun before asserting performance of a later
revision or a different experiment.}
\begin{tabularx}{\linewidth}{@{}lY@{}}\toprule
Source & Evidence carried into this manuscript\\\midrule
\code{docs/native-kernels.md} & First native/preceding-FP32 comparisons and numerical boundaries\\
\code{docs/native-kernels-v2.md} & Native-to-native 180-step measurements, matched hashes and data-path fixtures\\
\code{docs/performance.md} & Earlier Python/save-path measurements; memory and operation-count scope\\
\code{docs/recording.md} & Exact-byte 401-frame replay representation comparison\\
\code{docs/running.md} & Compact-checkpoint size fixture and continuation semantics\\
\code{docs/binary-storage.md} & Deep-model archive fixture and lossless storage rules\\
\code{docs/releases/0.0.0.24.md} & Verified 20-run A2C binary conversion totals\\\bottomrule
\end{tabularx}
\end{table}

Source commits for the historical native stages are \code{baa5101} and
\code{abd310d}; storage/conversion and champion changes were introduced in their
own documented rollouts. These identifiers locate code stages, not a guarantee
that ignored raw measurement artifacts exist in every clone. Preserve raw reports
alongside the eventual public experiment release.
